% ============================================================
%  CHAPITRE II : État de l'art
% ============================================================
\chapter{État de l'art}
\label{chap:etat_art}

\section*{Introduction}

Ce chapitre présente une revue de la littérature couvrant les domaines de l'\gls{ia} et de la cybersécurité. Nous y analysons les approches existantes, les architectures de modèles utilisés et les métriques d'évaluation pertinentes.

% -------------------------------------------------------
\section{Intelligence artificielle et apprentissage automatique}
\label{sec:ia_ml}

\subsection{Définitions fondamentales}

\iacsdefinition{Machine Learning}{
  Le \gls{ml} est une branche de l'\gls{ia} qui permet aux machines d'apprendre à partir de données sans être explicitement programmées. On distingue trois paradigmes principaux : l'apprentissage supervisé, non supervisé et par renforcement (\gls{rl}).
}

\subsection{Réseaux de neurones convolutifs}

Les \gls{cnn} sont des architectures de \gls{dl} particulièrement adaptées au traitement d'images et de données structurées. L'opération de convolution est définie par l'équation~\eqref{eq:convolution} :

\begin{equation}
  (f * g)(t) = \int_{-\infty}^{+\infty} f(\tau) \cdot g(t - \tau) \, d\tau
  \label{eq:convolution}
\end{equation}

La sortie d'une couche de convolution discrète à deux dimensions s'exprime comme :

\begin{equation}
  S(i, j) = \sum_{m} \sum_{n} I(i+m, j+n) \cdot K(m, n) + b
  \label{eq:conv2d}
\end{equation}

où $I$ est l'entrée, $K$ le noyau de convolution et $b$ le biais.

\subsection{Fonctions d'activation}

Le tableau~\ref{tab:activations} présente les principales fonctions d'activation utilisées dans les réseaux de neurones.

\begin{table}[H]
  \centering
  \caption{Fonctions d'activation courantes}
  \label{tab:activations}
  \renewcommand{\arraystretch}{1.5}
  \begin{tabular}{|l|c|c|}
    \hline
    \rowcolor{iacsTableHead}
    \textcolor{white}{\textbf{Fonction}} & 
    \textcolor{white}{\textbf{Expression}} & 
    \textcolor{white}{\textbf{Plage}} \\
    \hline
    ReLU & $f(x) = \max(0, x)$ & $[0, +\infty[$ \\
    \hline
    \rowcolor{iacsLightGray}
    Sigmoid & $\sigma(x) = \dfrac{1}{1 + e^{-x}}$ & $]0, 1[$ \\
    \hline
    Tanh & $\tanh(x) = \dfrac{e^x - e^{-x}}{e^x + e^{-x}}$ & $]-1, 1[$ \\
    \hline
    \rowcolor{iacsLightGray}
    Softmax & $\sigma(z_i) = \dfrac{e^{z_i}}{\sum_{j} e^{z_j}}$ & $]0, 1[$ \\
    \hline
  \end{tabular}
\end{table}

% -------------------------------------------------------
\section{Cybersécurité et détection d'intrusions}
\label{sec:cyber}

\subsection{Taxonomie des attaques réseau}

Les attaques réseau peuvent être classifiées en plusieurs catégories :

\begin{enumerate}
  \item \textbf{Attaques de déni de service (\gls{ddos})} : visant à rendre un service indisponible
  \item \textbf{Attaques d'injection (\gls{xss}, SQL Injection)} : exploitant les failles des applications web
  \item \textbf{Attaques de type Man-in-the-Middle} : interceptant les communications
  \item \textbf{Attaques par force brute} : tentant de deviner les identifiants
\end{enumerate}

\subsection{Systèmes de détection d'intrusions}

Un \gls{ids} analyse le trafic réseau pour identifier les activités malveillantes. L'algorithme~\ref{algo:ids} décrit le fonctionnement simplifié d'un \gls{ids} basé sur le \gls{ml}.

\begin{algorithm}[H]
  \caption{Détection d'intrusions par apprentissage supervisé}
  \label{algo:ids}
  \KwEntree{Flux réseau $\mathcal{F} = \{f_1, f_2, \ldots, f_n\}$, Modèle entraîné $\mathcal{M}$}
  \KwSortie{Étiquettes de classification $\mathcal{Y} = \{y_1, y_2, \ldots, y_n\}$}

  \Pour{chaque flux $f_i \in \mathcal{F}$}{
    $\mathbf{x}_i \leftarrow$ \texttt{extraireCaractéristiques}($f_i$)\;
    $\hat{y}_i \leftarrow \mathcal{M}(\mathbf{x}_i)$\;
    \uSi{$\hat{y}_i = \text{«\,malveillant\,»}$}{
      \texttt{générerAlerte}($f_i$, $\hat{y}_i$)\;
      \texttt{loguer}($f_i$, $\hat{y}_i$, horodatage)\;
    }
    \Sinon{
      \texttt{autoriser}($f_i$)\;
    }
  }
  \Retour{$\mathcal{Y}$}
\end{algorithm}

% -------------------------------------------------------
\section{Métriques d'évaluation}
\label{sec:metriques}

Les performances des modèles sont évaluées à l'aide des métriques suivantes :

\begin{align}
  \text{Précision}   &= \frac{VP}{VP + FP} \label{eq:precision} \\[6pt]
  \text{Rappel}      &= \frac{VP}{VP + FN} \label{eq:rappel} \\[6pt]
  \text{F1-score}    &= 2 \cdot \frac{\text{Précision} \times \text{Rappel}}{\text{Précision} + \text{Rappel}} \label{eq:f1}
\end{align}

où $VP$ désigne les vrais positifs, $FP$ les faux positifs et $FN$ les faux négatifs.

% -------------------------------------------------------
\section{Exemple de code source}
\label{sec:code}

Le listing~\ref{lst:python} montre un exemple de prétraitement de données réseau en Python.

\begin{lstlisting}[language=Python, caption={Prétraitement des données réseau}, label={lst:python}]
import pandas as pd
from sklearn.preprocessing import StandardScaler
from sklearn.model_selection import train_test_split

# Chargement du jeu de donnees
df = pd.read_csv("network_traffic.csv")

# Selection des caracteristiques
features = ['duration', 'src_bytes', 'dst_bytes',
            'num_failed_logins', 'count', 'srv_count']
X = df[features]
y = df['label']  # normal / attaque

# Normalisation
scaler = StandardScaler()
X_scaled = scaler.fit_transform(X)

# Separation train / test
X_train, X_test, y_train, y_test = train_test_split(
    X_scaled, y, test_size=0.2, random_state=42
)
\end{lstlisting}

% -------------------------------------------------------
\section{Étude comparative}
\label{sec:comparative}

Le tableau~\ref{tab:comparative} compare les travaux récents dans le domaine.

\begin{table}[H]
  \centering
  \caption{Étude comparative des approches existantes}
  \label{tab:comparative}
  \renewcommand{\arraystretch}{1.3}
  \small
  \begin{tabularx}{\textwidth}{|l|c|c|X|}
    \hline
    \rowcolor{iacsTableHead}
    \textcolor{white}{\textbf{Réf.}} & 
    \textcolor{white}{\textbf{Méthode}} & 
    \textcolor{white}{\textbf{Précision}} & 
    \textcolor{white}{\textbf{Observations}} \\
    \hline
    {[1]} & CNN & 94,5\% & Temps d'entraînement élevé \\
    \hline
    \rowcolor{iacsLightGray}
    {[2]} & LSTM & 93,2\% & Adapté aux séries temporelles \\
    \hline
    {[3]} & Random Forest & 91,8\% & Interprétabilité limitée \\
    \hline
    \rowcolor{iacsLightGray}
    {[4]} & Autoencoder & 92,1\% & Non supervisé \\
    \hline
    \textbf{Notre approche} & \textbf{CNN hybride} & \textbf{96,3\%} & \textbf{Temps réel} \\
    \hline
  \end{tabularx}
\end{table}

\section*{Conclusion}

Ce chapitre a permis de dresser un panorama des techniques existantes en matière de détection d'intrusions basée sur l'\gls{ia}. L'analyse comparative a mis en évidence les limites des approches actuelles, justifiant ainsi la solution proposée dans le chapitre suivant.
